% !TeX root = ../教材题库_全部条目_转写工作稿.tex
% Source: XJ-B1-C05-LAB-P198-Y0145
% 第5章 / 5.4 函数 y=A sin(ωx+φ) 的图象与性质 / 教材 p.198 / PDF 第205页
% OCR未校对：公式与文字均可能有误；未完成数学TeX转写；禁止直接发给学生
% 原题及相邻原解截图（核对用）：../题目原图/XJ-B1-C05-LAB-P198-Y0145-01.webp
% 本文件只在转写完成后作为题干使用；OCR 条目尚不是可靠题干。
\begin{bankitem}{用计算机作三角函数图象}{第5章 / 5.4 函数 y=A sin(\ensuremath{\omega}x+\ensuremath{\varphi}) 的图象与性质 / 教材 p.198 / PDF 第205页}{XJ-B1-C05-LAB-P198-Y0145}{OCR未校对：公式与文字均可能有误；未完成数学TeX转写；禁止直接发给学生}
\begin{ocrdraft}
[教材p.198 / PDF第205页]
\par
数学实验
\par
用计算机作函数3，-Asin（-+的图象
\par
借助“超级画板”或其他有类似功能的软件，可在计算机上作出函数y一
\par
Asin(o-\ocrglyph{00BB}x十@的动态图象．拖动参数A，仞和能观察到图象随参数变化的情形．
\par
操作步骤如下：
\par
1.输入函数解析式：在右键菜单里单击“函数或参数方程曲线”，在打开的对
\par
话框里选'3一丆0）”，在激活的输入栏里填写关sin(b关x十c),曲线的点数设置
\par
为1000，参数范围为一15到15，单击“确定”．
\par
2．作出可控制参数点：作坐标点，
\par
一2）
\par
一3），顺次标注为
\par
A,和自这三点分别向轴引垂线．
\par
3．作出cos的曲线，作为比较，线型可选择虚线．
\par
4．测量变量a，b，c，在测量值显示框里将等号左边的变量名分别改为A,和
\par
5．加上标题“正弦波的振幅、频率和相角”和"y=Asin()x十匆"，再添加上
\par
其他必要的说明\ensuremath{\cdot}
\par
6．用鼠标分别拖动点A，仞和或选择其中一点用左右箭头键驱动，观察图
\par
象的变化（如下图）．
\par
」4=丨．8042
\par
=2．2583
\par
巨0@0\ensuremath{\text{[OCR:GBP]}}0
\par
匣亘座到
\par
0
\par
图中虚线是3-cosx的图象．
\par
调整参数A，仞和使动态曲线和虚线尽可能重合．这时，三个参数的值分别
\par
是多少？
\end{ocrdraft}
\end{bankitem}
